\documentclass[12pt,a4paper]{article}
\usepackage[utf8]{inputenc}
\usepackage[T1]{fontenc}
\usepackage[spanish]{babel}
\usepackage[margin=2cm]{geometry}
\usepackage{xcolor}
\usepackage{amsmath}
\usepackage{enumitem}
\usepackage{microtype}
\setlength{\parindent}{0pt}
\setlength{\parskip}{0.8em}
\definecolor{azul}{HTML}{1E5A88}
\definecolor{fondo}{HTML}{F0F6FA}
\newcommand{\tiempo}[1]{\textcolor{azul}{\textbf{Tiempo: #1}}}
\newcommand{\voz}{\textcolor{azul}{\textbf{Narración -- Benjamin De La Fuente}}}
\newcommand{\clave}[1]{\fcolorbox{azul}{fondo}{\parbox{0.91\linewidth}{#1}}}
\begin{document}
\begin{center}
{\LARGE \textbf{Guion de narración}}\\[4pt]
{\large Benjamin De La Fuente -- Grupo 16}\[3pt]
Integral de velocidad, sumas de Riemann y desplazamiento
\end{center}
\hrule
\vspace{0.8em}
\clave{\textbf{Idea general:} Tu parte dura aproximadamente 1 minuto y 45 segundos, desde el minuto 1:45 hasta el 3:30 del video. Explicas la gráfica de velocidad, las sumas de Riemann y el cálculo del desplazamiento.}

\section*{1. Gráfica velocidad-tiempo}
\tiempo{1:45 -- 2:20}\par
\voz
\begin{quote}
Soy Benjamin De La Fuente. En esta gráfica, el eje horizontal es tiempo en segundos; el vertical es velocidad en metros por segundo, no posición.

El punto verde y el auto usan el mismo reloj. El tramo superior corresponde a la ida. El tramo sobre el eje indica la compra: el tiempo avanza y el auto sigue quieto.

El tramo inferior indica el regreso. Velocidad negativa significa cambiar de sentido, no perder distancia. Las líneas horizontales representan velocidades constantes.
\end{quote}
\textbf{Qué debes señalar al hablar:}
\begin{itemize}[leftmargin=1.4em]
\item Eje horizontal: tiempo $t$, medido en segundos.
\item Eje vertical: velocidad $v(t)$, medida en metros por segundo.
\item De $0$ a $15$ s: $v(t)=+1$ m/s; el auto va hacia la panadería.
\item De $15$ a $25$ s: $v(t)=0$; el auto está detenido durante la compra.
\item De $25$ a $40$ s: $v(t)=-1$ m/s; el auto vuelve hacia la casa.
\end{itemize}
\clave{\textbf{Cuidado:} La gráfica representa velocidad, no posición. El tramo inferior no dice que el auto esté bajo tierra; indica que se mueve en el sentido negativo de la recta.}

\section*{2. De sumas de Riemann a integral}
\tiempo{2:20 -- 2:55}\par
\voz
\begin{quote}
En cada intervalo pequeño, velocidad por duración aproxima un cambio de posición. Dividimos los cuarenta segundos en $n$ intervalos iguales, de ancho delta $t$, igual a cuarenta dividido por $n$.

Elegimos el punto medio, $t_i^*$. Sigma significa sumar esos aportes. Mostramos ocho, dieciséis y treinta y dos intervalos.

Al refinar la partición, delta $t$ tiende a cero y obtenemos la integral. Aquí los cortes coinciden con las velocidades constantes, por eso las sumas ya son exactas.
\end{quote}
\textbf{Fórmula que aparece:}
\[
\sum_{i=1}^{n} v(t_i^*)\,\Delta t \longrightarrow \int_0^{40} v(t)\,dt.
\]
\textbf{Cómo leerla en voz alta:}

``La suma, desde $i$ igual a uno hasta $n$, de velocidad evaluada en el punto medio $t$ sub $i$ estrella, multiplicada por delta $t$, tiende a la integral de cero a cuarenta de velocidad respecto del tiempo.''

\textbf{Símbolos importantes:}
\begin{itemize}[leftmargin=1.4em]
\item $n$: cantidad de subintervalos.
\item $\Delta t=40/n$: ancho de cada subintervalo, en segundos.
\item $t_i^*$: punto medio del subintervalo número $i$.
\item $\Sigma$: suma de todos los aportes pequeños.
\item $dt$: integración respecto del tiempo.
\end{itemize}
\clave{En este ejemplo, las fronteras de la partición coinciden con $t=15$ y $t=25$. Por eso, aunque usemos un número finito de rectángulos, cada rectángulo tiene una velocidad constante y la suma es exacta. Esto no ocurre necesariamente cuando la velocidad varía de forma continua.}

\section*{3. Integral de velocidad: desplazamiento}
\tiempo{2:55 -- 3:30}\par
\voz
\begin{quote}
En el primer tramo, un metro por segundo durante quince segundos aporta más quince metros. La compra aporta cero por diez: cero metros.

En el regreso, menos un metro por segundo durante quince segundos aporta menos quince metros. Cada región tiene área geométrica positiva de quince metros, pero la región inferior contribuye con signo negativo.

La integral entre cero y cuarenta suma esos aportes: quince más cero menos quince da cero metros de desplazamiento. Las unidades son metros por segundo multiplicados por segundos.
\end{quote}
\textbf{Cálculo que aparece:}
\begin{align*}
\Delta x &= \int_0^{15} 1\,dt + \int_{15}^{25} 0\,dt + \int_{25}^{40} (-1)\,dt \\
&= 15+0-15 \\
&= 0\ \mathrm{m}.
\end{align*}
\textbf{Interpretación:}
\begin{itemize}[leftmargin=1.4em]
\item Ida: $+1$ m/s $\times$ $15$ s $=+15$ m.
\item Compra: $0$ m/s $\times$ $10$ s $=0$ m.
\item Regreso: $-1$ m/s $\times$ $15$ s $=-15$ m.
\item Resultado: $+15+0-15=0$ m.
\end{itemize}
\clave{\textbf{Conclusión de tu bloque:} Integrar la velocidad calcula desplazamiento, porque conserva el signo. Los $15$ metros de ida y los $15$ metros de regreso se cancelan: el auto termina donde comenzó. Esto no significa que haya recorrido cero metros; la distancia se explicará en el bloque siguiente.}

\section*{Repaso rápido antes de grabar}
\begin{enumerate}[leftmargin=1.4em]
\item La velocidad negativa indica dirección opuesta, no distancia negativa.
\item El área bajo el eje de velocidad cuenta negativamente para el desplazamiento.
\item La unidad final es metro: $(\mathrm{m/s})\cdot\mathrm{s}=\mathrm{m}$.
\item $\int v(t)\,dt$ entrega desplazamiento; $\int |v(t)|\,dt$ entrega distancia.
\item En este caso, $x(40)-x(0)=0-0=0$ m confirma el resultado.
\end{enumerate}
\end{document}
